\documentclass[a4paper,10pt]{article}
\usepackage{lamia-pedagogy}

\begin{document}

% =========================================================================
% PAGE 1 : Fiche Enseignant & Préparation à l'Écrit
% =========================================================================
\ficheHeader{1 Heure}{\makebox[2.5cm]{\dotfill}}{1\textsuperscript{ère} AM}{2026 / 2027}

\vspace{-3pt}
\competenceBanner{\faPen\ Projet 01 / Séquence 01 — Séance 08 : Préparation à l'Écrit}

\vspace{-2pt}
\modaliteSupportBar{Travail individuel guidé}{Situation d'intégration / Journal mural du collège}

\vspace{1pt}

% Consigne d'écriture
\begin{tcolorbox}[
  enhanced,
  colback=accentlight!30,
  colframe=accent,
  arc=2mm,boxrule=0.8pt,
  top=4pt,bottom=4pt,left=8pt,right=8pt,
  title={\small\textbf{\textcolor{accent!90!black}{\faFileSignature\ Situation d'intégration}}},
  colbacktitle=white,coltitle=accent!90!black,
  attach boxed title to top left={yshift=-2mm,xshift=4mm},
  boxed title style={boxrule=0.6pt,colframe=accent,arc=1.2mm,top=1pt,bottom=1pt}
]
\vspace{2pt}
{\small
\textbf{Consigne :} Dans le cadre de la semaine de sensibilisation à la santé au collège, rédige un court texte explicatif (de 4 à 6 lignes) pour expliquer à tes camarades \textbf{l'importance d'une bonne hygiène corporelle}.
}
\end{tcolorbox}

\vspace{3pt}

% Entraînement préparatoire
\begin{tcolorbox}[
  enhanced,
  colback=white,
  colframe=cardborder,
  arc=2mm,boxrule=0.8pt,
  top=4pt,bottom=4pt,left=8pt,right=8pt,
  title={\small\textbf{\textcolor{primary}{\faLayerGroup\ Activités préparatoires : Remets les mots en ordre}} \hfill \badge[primary]{Entraînement}},
  coltitle=primary,colbacktitle=primarylight,titlerule=0.5pt,titlerule style=cardborder
]
{\small
Reconstitue les phrases explicatives en remettant les mots dans l'ordre :
\begin{enumerate}[label=\textbf{\arabic*.},leftmargin=*,noitemsep,topsep=2pt]
  \item \textit{des saletés / consiste à / Le lavage des mains / se débarrasser .}\\
    $\longrightarrow$ \drule{\linewidth}
  \item \textit{est un ensemble / L'hygiène du corps / de propreté quotidienne / de règles .}\\
    $\longrightarrow$ \drule{\linewidth}
  \item \textit{sont / minuscules / des êtres vivants / Les microbes .}\\
    $\longrightarrow$ \drule{\linewidth}
\end{enumerate}
}
\end{tcolorbox}

\vspace{3pt}

% Grille des critères de réussite
\begin{tcolorbox}[
  enhanced,
  colback=white,
  colframe=secondary,
  arc=2mm,boxrule=0.8pt,
  top=4pt,bottom=4pt,left=7pt,right=7pt,
  title={\small\textbf{\textcolor{secondary!90!black}{\faClipboardCheck\ Critères de réussite (Pour réussir ma rédaction)}}},
  colbacktitle=secondarylight,coltitle=secondary!90!black,titlerule=0.5pt
]
{\footnotesize
\begin{tabularx}{\linewidth}{|l|X|c|}
\hline
\rowcolor{primarylight}\textbf{Composantes} & \textbf{Critères à respecter} & \textbf{Vérifié} \\
\hline
1. Titre & Donner un titre adapté au sujet (\textit{ex: L'importance de l'hygiène}). & \qcmCircle \\
\hline
2. Structure & Respecter les 3 phases : Questionnement, Explication, Bilan. & \qcmCircle \\
\hline
3. Vocabulaire & Employer un procédé explicatif (Définition ou Synonyme). & \qcmCircle \\
\hline
4. Grammaire & Utiliser une proposition subordonnée relative par \textbf{\guillemotleft\ qui \guillemotright}. & \qcmCircle \\
\hline
5. Conjugaison & Conjuguer les verbes d'action au \textbf{présent de l'indicatif}. & \qcmCircle \\
\hline
\end{tabularx}
}
\end{tcolorbox}

\newpage

% =========================================================================
% PAGE 2 : Fiche Élève (Rédaction du 1er Jet)
% =========================================================================
\ficheHeader{1 Heure}{Nom : \makebox[4cm]{\dotfill}}{1\textsuperscript{ère} AM}{Classe : \makebox[1.5cm]{\dotfill}}

\vspace{-3pt}
\competenceBanner{\faPenFancy\ Fiche de Rédaction : Mon Texte Explicatif (1\textsuperscript{er} Jet)}

\vspace{2pt}

% Boîte à mots d'aide
\begin{tcolorbox}[
  enhanced,
  colback=secondarylight!30,
  colframe=cardborder,
  arc=1.5mm,boxrule=0.6pt,
  top=3pt,bottom=3pt,left=6pt,right=6pt,
  title={\scriptsize\textbf{\textcolor{primary}{\faTools\ Boîte à mots d'aide pour la rédaction}}},
  coltitle=primary,colbacktitle=white,titlerule=0.4pt
]
{\scriptsize
\begin{tabularx}{\linewidth}{X X X}
\textbf{Noms :} l'hygiène, la propreté, le savon, les microbes, la santé, la douche, les mains. &
\textbf{Verbes :} consister à, être, éliminer, laver, frotter, protéger, éviter, préserver. &
\textbf{Connecteurs / Mots outils :} d'abord, puis, qui, parce que, car, donc, enfin. \\
\end{tabularx}
}
\end{tcolorbox}

\vspace{3pt}

% Espace d'écriture de l'élève
\begin{tcolorbox}[
  enhanced,
  colback=white,
  colframe=primary,
  arc=2mm,boxrule=0.8pt,
  top=6pt,bottom=6pt,left=10pt,right=10pt,
  title={\small\textbf{\textcolor{white}{\faPen\ Espace de rédaction de l'élève}}},
  coltitle=white,colbacktitle=primary
]
\vspace{2pt}
\begin{center}
\textbf{Titre :} \drule{9cm}
\end{center}

\vspace{8pt}
\textbf{Phase de questionnement (Introduction) :}\\
\drule{\linewidth}\\[10pt]
\drule{\linewidth}

\vspace{10pt}
\textbf{Phase d'explication (Développement) :}\\
\drule{\linewidth}\\[10pt]
\drule{\linewidth}\\[10pt]
\drule{\linewidth}\\[10pt]
\drule{\linewidth}

\vspace{10pt}
\textbf{Phase de conclusion :}\\
\drule{\linewidth}\\[10pt]
\drule{\linewidth}
\vspace{4pt}
\end{tcolorbox}

\vspace{3pt}

% Grille d'auto-évaluation finale
\begin{tcolorbox}[
  enhanced,
  colback=neutrallight,
  colframe=secondary!60!cardborder,
  arc=1.5mm,boxrule=0.6pt,
  top=3pt,bottom=3pt,left=7pt,right=7pt,
  title={\scriptsize\textbf{\faCheckDouble\ Grille de relecture de l'élève}},
  coltitle=primary,colbacktitle=secondarylight,titlerule=0.4pt
]
{\footnotesize
\begin{tabularx}{\linewidth}{Xcc}
\textbf{Critère de réussite :} & \textbf{Oui \faSmile} & \textbf{Non \faFrown} \\
\midrule
J'ai donné un titre en relation avec le thème de la propreté. & \qcmCircle & \qcmCircle \\
J'ai utilisé une relative par \guillemotleft\ qui \guillemotright\ et un verbe au présent. & \qcmCircle & \qcmCircle \\
J'ai soigné mon écriture et ma ponctuation. & \qcmCircle & \qcmCircle \\
\end{tabularx}
}
\end{tcolorbox}

\end{document}
